\section{Preliminares}

\begin{Theorem}{Teorema fundamental del calculo}
    - Sea $f \in \mathcal{C}^1$ en $\complex$. Entonces se tiene que:
    $$ f(b) - f(a) = \int_{\gamma[a, b]} f'(z) dz $$
\end{Theorem}

\begin{Lemma}{Lema del valor absoluto constante}
    - $ f : D \to \complex $ holomorfa: $|f(z)| = c \then f(z) = c_0 $
\end{Lemma}

\begin{Theorem}{Formula del valor medio}
    - $$ f(z_0) = \frac{1}{2 \pi} \int_{0}^{2\pi} f(z_0 + re^{it}) dt $$
\end{Theorem}

\begin{Proposition}{Sobre la funcion constante}
    - $f $ continua es constante $\iff max = min = c $ 
\end{Proposition}

\begin{Theorem}{Principio del máximo}
    - $ f: D \to \complex $ holomorfa y $max_{z \in D} |f(z)| \leq |f(z_0)|$
    entonces $f(z)$ es constante.
\end{Theorem}

\begin{Theorem}{Teorema fundamental del algebra (debil)}
    - Sea $P_N(z) $ un polinomio de grado N $\then \exists z_0 \in D: P_N(z_0) = 0$
\end{Theorem}